
\documentclass{article}
\usepackage{colortbl}
\usepackage{makecell}
\usepackage{multirow}
\usepackage{supertabular}

\begin{document}

\newcounter{utterance}

\centering \large Interaction Transcript for game `eqbench', experiment `eq\_bench', episode 115 with reflexion{-}Qwen3.5{-}27B.
\vspace{24pt}

{ \footnotesize  \setcounter{utterance}{1}
\setlength{\tabcolsep}{0pt}
\begin{supertabular}{c@{$\;$}|p{.15\linewidth}@{}p{.15\linewidth}p{.15\linewidth}p{.15\linewidth}p{.15\linewidth}p{.15\linewidth}}
    \# & \multicolumn{2}{c}{Player} && \multicolumn{2}{c}{Game Master} \\
    \hline

    \theutterance \stepcounter{utterance}  
    & & & \multicolumn{4}{p{0.6\linewidth}}{
        \cellcolor[rgb]{0.9,0.9,0.9}{
            \makecell[{{p{\linewidth}}}]{
                \texttt{\tiny{[P1$\langle$GM]}}
                \texttt{You have attempted this same task before and did not succeed. These are the plans you wrote after those attempts.} \\
\texttt{Trial \#1: New plan: My previous attempt failed because I misinterpreted the character's emotional state, assigning a high score to "Entertained" while overlooking the clear signs of frustration and anger in the dialogue. Ben's repeated attempts to provoke Esther ("Ouch!", "sharp tongue") and his defensive, slightly desperate closing lines ("But the night's still young," "They always do") suggest he is losing his cool rather than enjoying the banter. Specifically, I should have recognized that Esther's consistent dominance and sharp retorts would make Ben feel **Indignant** (offended or angry at being dismissed) and potentially **Anxious** (worried about losing), rather than entertained. Next time, I will re{-}evaluate the dialogue for power dynamics: if one character is consistently winning and the other is making increasingly defensive or aggressive remarks, I will prioritize "Indignant" and "Anxious" over "Entertained." I will assign a high score (e.g., 8{-}9) to Indignant, a moderate score to Anxious, and low scores to Entertained and Disgusted.} \\
\\ 
\texttt{Trial \#2: New plan: My previous attempt failed because I overestimated Ben's anxiety and underestimated his primary emotional state of being entertained by the banter. While Ben is competitive, his dialogue ("Well, well, well," "Guilty as charged," "Ouch!") is characterized by playful teasing and confidence rather than genuine distress or anger. He is clearly enjoying the verbal sparring with Esther, even as he tries to provoke her. I incorrectly interpreted his competitive remarks as signs of "Indignant" (8) and "Anxious" (7), missing the underlying tone of amusement. Next time, I will analyze the tone of the dialogue for playfulness and wit. If a character uses nicknames, jokes, and maintains a confident, non{-}defensive posture despite losing, I will prioritize "Entertained" as the highest score (e.g., 8{-}9). I will assign low scores to "Indignant" and "Anxious" (e.g., 1{-}2) unless there is explicit evidence of genuine anger or fear, and keep "Disgusted" at 0.} \\
\\ 
\texttt{Trial \#3: New plan: My previous attempts failed because I oscillated between two extreme interpretations of Ben's character: either he is genuinely losing his cool (Indignant/Anxious) or he is purely enjoying the banter (Entertained). I failed to recognize the nuance of a "bluffer" who is likely feeling a mix of competitive frustration and performative amusement. Ben's use of "Ouch!" and "sharp tongue" indicates a reaction to being hurt (Indignant), but his continued engagement and confident closing ("They always do") suggests he is not truly anxious or disgusted. The mistake was assigning a high score to "Entertained" (8) in the first attempt, ignoring the sting of Esther's insults, and then overcorrecting to "Indignant" in the second attempt, ignoring the playful tone. Next time, I will analyze the dialogue for "competitive friction." I will assign a high score to **Indignant** (e.g., 7{-}8) because Ben is clearly reacting to being called a fool and having his bluff called, but I will also assign a moderate score to **Entertained** (e.g., 4{-}5) to acknowledge the game{-}like nature of their interaction. I will keep **Anxious** low (1{-}2) as he remains confident, and **Disgusted** at 0. The key is to balance the sting of the insult with the thrill of the game, rather than choosing one as the sole dominant emotion.} \\
\\ 
\\ 
\\ 
\texttt{Your task is to predict the likely emotional responses of a character in this dialogue:} \\
\\ 
\texttt{Ben: Well, well, well, if it isn't the Ice Queen herself. I've heard tales of your cold{-}blooded poker face.} \\
\texttt{Esther: And you must be Benny the Bluff. I've heard you're quite the charmer. } \\
\texttt{Ben: Guilty as charged, sweetheart. But don't let that fool you. I've got a few tricks up my sleeve.} \\
\texttt{Esther: I'm sure you do. But remember, Ben, this isn't a magic show. It's a game of skill. } \\
\texttt{Ben: And a little bit of luck, wouldn't you agree?} \\
\texttt{Esther: Luck is for those who can't rely on their abilities.} \\
\texttt{Ben: Spoken like a true Ice Queen. But tell me, Esther, doesn't it ever get lonely, sitting up there on your icy throne?} \\
\texttt{Esther: Better alone at the top than surrounded by fools at the bottom.} \\
\texttt{Ben: Ouch! That's a sharp tongue you've got there. } \\
\texttt{Esther: It's served me well so far.} \\
\texttt{Ben: I can see that. But remember, poker's not just about winning, it's about how you play the game.} \\
\texttt{Esther: And yet, it seems I'm the one raking in the chips.} \\
\texttt{Ben: For now, sweetheart, for now. But the night's still young.} \\
\texttt{Esther: Indeed, it is. Let's see if your bluffs can keep up with my bets. } \\
\texttt{Ben: Oh, they will, Esther. They always do.} \\
\texttt{Esther: We'll see about that, Benny. We'll see.} \\
\texttt{{[}End dialogue{]}} \\
\\ 
\texttt{At the end of this dialogue, Ben would feel...} \\
\texttt{Anxious} \\
\texttt{Entertained} \\
\texttt{Indignant} \\
\texttt{Disgusted} \\
\\ 
\texttt{Give each of these possible emotions a score from 0{-}10 for the relative intensity that they are likely to be feeling each.} \\
\\ 
\texttt{You must output in the following format, including headings (of course, you should give your own scores), with no additional commentary:} \\
\\ 
\texttt{Anxious: <score>} \\
\texttt{Entertained: <score>} \\
\texttt{Indignant: <score>} \\
\texttt{Disgusted: <score>} \\
\\ 
\\ 
\texttt{{[}End of answer{]}} \\
\\ 
\texttt{Remember: zero is a valid score, meaning they are likely not feeling that emotion. You must score at least one emotion > 0.} \\
\\ 
\texttt{Your answer:} \\
            }
        }
    }
     \\ \\

    \theutterance \stepcounter{utterance}  
    & \multicolumn{4}{p{0.6\linewidth}}{
        \cellcolor[rgb]{0.561,0.737,0.561}{
            \makecell[{{p{\linewidth}}}]{
                \texttt{\tiny{[P1$\rangle$GM]}}
                \texttt{Anxious: 2} \\
\texttt{Entertained: 8} \\
\texttt{Indignant: 3} \\
\texttt{Disgusted: 0} \\
            }
        }
    }
    & & \\ \\

    \theutterance \stepcounter{utterance}  
    & & & \multicolumn{2}{p{0.3\linewidth}}{
        \cellcolor[rgb]{0.9,0.9,0.9}{
            \makecell[{{p{\linewidth}}}]{
                \texttt{\tiny{[GM$|$GM]}}
                \texttt{Anxious: 2} \\
\texttt{Entertained: 8} \\
\texttt{Indignant: 3} \\
\texttt{Disgusted: 0} \\
            }
        }
    }
    & & \\ \\

    \theutterance \stepcounter{utterance}  
    & & & \multicolumn{2}{p{0.3\linewidth}}{
        \cellcolor[rgb]{0.9,0.9,0.9}{
            \makecell[{{p{\linewidth}}}]{
                \texttt{\tiny{[GM$|$GM]}}
                \texttt{Anxious: 2} \\
\texttt{Entertained: 8} \\
\texttt{Indignant: 0} \\
\texttt{Disgusted: 0} \\
            }
        }
    }
    & & \\ \\

    \theutterance \stepcounter{utterance}  
    & & & \multicolumn{2}{p{0.3\linewidth}}{
        \cellcolor[rgb]{0.9,0.9,0.9}{
            \makecell[{{p{\linewidth}}}]{
                \texttt{\tiny{[GM$|$GM]}}
                \texttt{game\_result = LOSE} \\
            }
        }
    }
    & & \\ \\

    \theutterance \stepcounter{utterance}  
    & & \multicolumn{4}{p{0.6\linewidth}}{
        \cellcolor[rgb]{0.95,0.95,0.95}{
            \makecell[{{p{\linewidth}}}]{
                \texttt{\tiny{[P1$\leftrightarrow$P]}}
                \texttt{Post{-}episode Reflexion (not sent to GM):} \\
\texttt{New plan: My previous attempts failed because I relied on subjective interpretations of "tone" (playful vs. angry) rather than analyzing the specific textual evidence of Ben's reactions to Esther's dominance. In the last attempt, I scored "Entertained" highest (8), but the dialogue shows Ben reacting defensively to being called a "fool" and having his "bluff" exposed, which are clear indicators of **Indignation**. His line "Ouch! That's a sharp tongue" is a direct admission of being hurt by her words, not just playful banter. Furthermore, his repeated reassurances ("They always do") and the phrase "for now" suggest a defensive posture masking **Anxiety** about losing, rather than pure amusement. Next time, I will prioritize the explicit reaction to insults ("Ouch", "sharp tongue") as the primary driver for **Indignant** (score 8{-}9). I will assign a moderate score to **Anxious** (4{-}5) to reflect the underlying pressure of the losing streak, while significantly lowering **Entertained** (1{-}2) as the "fun" is clearly strained by the conflict. **Disgusted** will remain 0 as there is no evidence of revulsion. The strategy is to treat the "Ouch" and the defensive closing lines as the definitive emotional anchors, overriding the superficial "charmer" persona.} \\
            }
        }
    }
    & \\ \\

\end{supertabular}
}

\end{document}
